\section{Representations directed into a prescribed cube}
\label{sec:maximum-cube-inward}

A maximum parent can supply more than an unspecified representation of
its marked section. It can supply a representation whose edges enter any
chosen contained cube. This resolves ordinary-map compatibility for
amalgams over a full common cube when all but at most one piece have
maximum parents. The argument uses the literal extension theorem of
Chalopin--Chepoi--Moran--Warmuth; it does not construct a maximum parent
of the amalgam.

Write $Q_S(a)=\{a^T:T\subseteq S\}$, where $a^T$ flips the coordinates
of $T$ in $a$. All coordinate sets are finite. A representation map of
an ample class $H$ is a bijection $r:H\to\operatorname{Sh}(H)$ such that
\[
 (u\mathbin\triangle v)\cap(r(u)\cup r(v))\ne\varnothing
 \qquad(u\ne v).
\]
Here $u\mathbin\triangle v$ denotes the coordinates on which the words
differ. We use the equivalent orientation description
\cite[Theorem~6.8]{ChalopiChepoiMoran22}: $r(v)$ is the set of outgoing
edge directions, the outgoing cube $Q_{r(v)}(v)$ is contained in $H$,
and every contained cube has exactly one sink. Such orientations are
called ordinary unique-sink orientations (USOs); directed cycles are allowed. The boundary at a
contained cube $Q_S(a)$ is \emph{inward} if every edge between that cube
and its complement points into the cube, equivalently
$r(v)\subseteq S$ for every $v\in Q_S(a)$.

\begin{theorem}[Inward representations from a maximum parent]
\label{thm:maximum-cube-inward}
Let $M\subseteq\{0,1\}^F$ be a nonempty maximum class and
$Q_S(a)\subseteq M$. There is a representation map $r$ of $M$ with
\begin{equation}
 r(a^T)=T\qquad(T\subseteq S).
 \label{eq:maximum-cube-inward-distance}
\end{equation}
Consequently its boundary at $Q_S(a)$ is inward.
Moreover, one can fix the entire orientation outside $Q_S(a)$, including
its boundary edges, so that every ordinary USO of that cube extends to
an ordinary USO of $M$.

The same conclusion about inward extension of every cube USO holds for
any nonempty coordinate section $C=M[\tau]$ and any cube contained in
$C$. The actual parent $M$, its rank, its auxiliary coordinates and its
mark $\tau$ can all be retained. Here $M[\tau]$ means fixing an
auxiliary coordinate set to $\tau$ and deleting those fixed coordinates.
The representation of $M$ may be chosen anew.
\end{theorem}

The singleton-cube case is the established prescribed-root refinement of
the maximum-class construction. The additional requirement here concerns
all vertices of a positive-dimensional cube simultaneously. The proof
selects its free directions first in the CCMW recursion. No corner
order, geometric realization, or acyclic exterior is assumed.

\begin{proof}
We first recall the exact extension input
\cite[Proposition~5.2]{ChalopiChepoiMoran22}. If $D\subseteq P$ are
maximum classes on the same coordinate set, of ranks $d-1$ and $d$,
then every representation map of $D$ extends literally to one of $P$.
The extension assigns each word of $P\setminus D$ the support of its
unique source cube. In particular, no label on $D$ changes.

Here is the associated lifting step, including the full-rank boundary
case. Fix a coordinate $e$, let $P=\pi_e M$ be the projection, and let
$D=M^e$ be the reduction consisting of projected two-point fibers.
If $M$ has rank $d$ with $1\le d<|F|$, the maximum reduction theorem
gives ranks $d$ and $d-1$ for $P$ and $D$. Equivalently, their two
Sauer bounds add to $|M|$ by Pascal's identity, so both bounds must be
attained. If $d=|F|$, then $M$ is the full cube and $P=D$ is the full
remaining cube.

Given a map $q$ of $D$, extend it literally to a map $p$ of $P$ by
Proposition~5.2, or set $p=q$ when $P=D$. For either bit $b$, define
\begin{equation}
 r(v)=p(\pi_e v)\cup
 \begin{cases}
  \{e\},&\pi_e v\in D\text{ and }v_e\ne b,\\
  \varnothing,&\text{otherwise}.
 \end{cases}
 \label{eq:maximum-cube-inward-lift}
\end{equation}
This is the lift in the proof of
\cite[Theorem~5.1]{ChalopiChepoiMoran22}, allowing either shore as the
sink of every two-point fiber. For completeness, the labels avoiding
$e$ occur once each: use the singleton lift over $P\setminus D$ and
the $b$-lift over $D$. The labels containing $e$ occur once each on
the opposite lifts over $D$. These are exactly the shattered supports
of $M$. If two words have different projections, nonclashing of $p$
separates them on a coordinate other than $e$. Otherwise they are the
two lifts of a point of $D$, and one label contains $e$. Thus the lift
is a representation map, also when $P=D$.

We prove \eqref{eq:maximum-cube-inward-distance} by induction on $|F|$,
handling the case $S=\varnothing$ within the same induction. If $M$
has rank zero, it is the singleton $\{a\}$ and the empty label works.
Otherwise, if $S=\varnothing$, choose an incident edge at $a$ in the
connected one-inclusion graph of the nontrivial maximum class $M$;
write $e$ for its direction. Then $\pi_e a\in D$. Induction supplies
a map $q$ of $D$ with $q(\pi_e a)=\varnothing$. Take $b=a_e$ in
\eqref{eq:maximum-cube-inward-lift}; its label at $a$ is empty.
This proves the prescribed-root case, including full cubes.

If $S\ne\varnothing$, choose $e\in S$. Both lifts of every word of
$Q_{S\setminus\{e\}}(\pi_e a)$ belong to $M$, so this projected cube
is contained in $D$. Induction gives a map $q$ of $D$ with
\[
 q((\pi_e a)^T)=T\qquad(T\subseteq S\setminus\{e\}).
\]
Extend to $P$ and lift with $b=a_e$. At $a^T$ the old label is
$T\setminus\{e\}$, and the new coordinate is added exactly when
$e\in T$. This proves the required formula for every $T\subseteq S$.

We next justify replacing the cube orientation while keeping one fixed
exterior. At a cube vertex every outgoing edge remains inside the cube,
so any replacement USO supplies its outgoing cube. At outside vertices
all outgoing sets remain unchanged. For the unique-sink condition,
let $B\subseteq M$ be a contained cube meeting $Q_S(a)$. Their
intersection $L$ is a face of both cubes. Its old unique sink is also
a sink of $B$, since all edges from $B\setminus L$ into $L$ point
inward. Thus every word of $B\setminus L$ had an outgoing edge within
$B$ and still does after replacement. The new orientation on $L$ has
one sink, also a sink of $B$. Cubes disjoint from $Q_S(a)$ are
unchanged. This proves the replacement assertion.

Finally, lift a cube $Q\subseteq C=M[\tau]$ to the cube with the same
free coordinates and auxiliary values $\tau$ in $M$. Apply the proved
maximum-class statement to that cube and restrict the resulting
orientation to the marked coordinate face. Restriction of an ordinary
USO to a coordinate-face intersection preserves both the outgoing-cube
and unique-sink conditions: outgoing fixed directions are removed,
and every cube of the section was a cube of the parent. Its boundary
at $Q$ remains inward. The lifted cube and its ordinary orientations
are naturally identical to $Q$, so every prescribed cube USO extends
with one fixed section exterior. No parent or mark has changed.
\end{proof}

\begin{theorem}[Common-cube amalgams with one prescribed piece]
\label{thm:maxmin-common-cube-map-selection}
Let $S,F_1,\ldots,F_m$ be pairwise disjoint coordinate sets, and let
$A_i\subseteq\{0,1\}^{S\sqcup F_i}$ be nonempty ample classes containing
$Q_S\times\{0\}^{F_i}$. Embed each piece with all other exclusive
blocks zero and put $H=\bigcup_i A_i$.
If all pieces except possibly $A_j$ belong to $\mathsf{MaxMin}$ and
$A_j$ has a supplied ordinary representation, then $H$ is ample and
has an ordinary representation preserving every supplied out-set of
$A_j$.

If all pieces belong to $\mathsf{MaxMin}$, every ordinary USO on the
common cube extends to $H$, with inward boundary there. One can keep
the same exterior and boundary for all choices of this cube USO.
\end{theorem}

\begin{proof}
Any support meeting two exclusive blocks is unshattered, because its
pattern positive in both blocks is absent. On $S\sqcup F_i$ the other
pieces add only patterns already realized by the common cube in $A_i$.
Consequently
\[
 \operatorname{Sh}(H)=\bigcup_i\operatorname{Sh}(A_i),\qquad
 \operatorname{Sh}(A_i)\cap\operatorname{Sh}(A_k)=2^S\quad(i\ne k).
\]
The pieces intersect pairwise in precisely $Q_S$, so
$|H|=\sum_i|A_i|-(m-1)2^{|S|}=|\operatorname{Sh}(H)|$.
The Sandwich equality proves ampleness.

Restrict the supplied orientation of $A_j$ to the common cube, obtaining
an ordinary cube USO $q$. For each $i\ne j$,
Theorem~\ref{thm:maximum-cube-inward} supplies an orientation of $A_i$
with inward boundary and cube orientation $q$. Orient the union by
these agreeing edge orientations. Outside the common cube all neighbors
belong to one piece. At a common-cube word, every outgoing exclusive
direction comes from $A_j$, so its entire outgoing cube is precisely
the outgoing cube supplied by that piece. All out-sets on $A_j$ are
therefore preserved. Every contained cube belongs to a piece: varying
two exclusive blocks would require a missing word positive in both,
and a fixed nonzero exclusive coordinate forces the cube into its
piece. Thus every contained cube has its unique sink in that piece.
These are the two ordinary-USO conditions.

When all pieces are maximum-generated, apply
Theorem~\ref{thm:maximum-cube-inward} inwardly to every piece with the
same prescribed cube orientation. Their fixed exteriors are independent
of that choice, and the preceding gluing proof applies without any
outgoing exclusive direction at the common cube. This proves the final
assertion.
\end{proof}

For example, if a hypothetical ample forbidden pc-minor of
$\mathsf{MaxMin}$ has such a decomposition with at least two nontrivial
pieces, each piece is a proper coordinate-face minor and hence belongs
to $\mathsf{MaxMin}$. The theorem represents this forbidden class.
It does not exclude that class from the forbidden basis: ordinary
representability is weaker than supplying an actual maximum parent.
For a forbidden dismantlable-ample input, the same conclusion needs
maximum ancestry of its pieces as an additional hypothesis; positivity
of a proper minor alone does not provide it.

\paragraph{A cyclic exterior is necessary in general.}
The cornerless maximum class on twelve coordinates with $299$ words
contains a three-cube. The constructor in
\path{../maxmin-forbidden/research/maximum_cube_map_constructor.py} gives the representation
in Theorem~\ref{thm:maximum-cube-inward} directly from its domain.
With the monotone cube orientation its exterior is necessarily cyclic:
a globally acyclic USO of a nonempty ample class has a source, whose
outgoing cube would make it a corner. An inward boundary prevents a
cycle from crossing the cube boundary, so with a monotone cube every
cycle lies outside it. Thus ordinary inward extension is strictly
more general than relative corner peeling to that cube.
The solver-free replay checks this map, a $129$-word marked section of
the same parent, replacement by a cyclic three-cube USO, and $555$- and
$590$-word common-cube unions. These last unions certify map gluing;
no forbiddenness or maximum ancestry of those unions is asserted.

\subsection{Noncube coordinate faces and prescribed section maps}
\label{sec:maximum-noncube-face-inward}

The full-cube hypothesis can be removed when the retained set is an
intersection with an ambient coordinate face. Two observations suffice.
First, the CCMW lift can direct every edge in one chosen coordinate
toward either shore. Second, replacement with inward boundary works for
an arbitrary coordinate-face intersection, since its intersection with
any contained cube is still a full face of that cube. The latter is
also the case $E=T=\varnothing$ of the
\href{../dismantlable-ample/sections/classification/coupled_face_replacement.tex}
{face-replacement rule with fixed exceptional out-sets}.
The additional step is to compose these facts along successive marked
sections; no compatibility of different supplied parent maps is required.

For $J\subseteq F$ and $\tau\in\{0,1\}^J$, write
\[
 H[\tau]=\{v\in H:v|_J=\tau\},
\]
retaining all original positions and constants in this subsection.
A map on $H[\tau]$ has no outgoing direction in $J$. Its boundary in
$H$ is inward when no edge from $H[\tau]$ to $H\setminus H[\tau]$
points outwards. The set $H[\tau]$ need not be a full cube.

\begin{theorem}[Prescribed maps on coordinate-face intersections]
\label{thm:maxmin-coordinate-face-inward}
Let $H=M[\alpha]$ be a nonempty marked section of a finite maximum
class $M$, and let $D=H[\tau]$ be a nonempty further coordinate-face
intersection. Every ordinary USO of $D$ extends to an ordinary USO
of $H$ with inward boundary at $D$. One fixed exterior, including
all boundary edges, works for every choice of the USO on $D$.

More strongly, for any chain of nonempty coordinate-face intersections
\[
 M=H_0\supseteq H_1\supseteq\cdots\supseteq H_k=D,
\]
every prescribed ordinary USO on $D$ extends to one on the same $M$
with inward boundary at every $H_i$. All out-sets outside $D$ can be
fixed independently of the prescribed map. In particular every
ordinary map on an actual marked section extends literally, with its
supports embedded in the original coordinate set, to its original
maximum parent. The parent, rank, auxiliaries and marks do not change.
\end{theorem}

\begin{proof}
We establish the single-coordinate supply before composing it. Suppose
$M$ has positive rank and choose any coordinate $e$ and either bit $b$.
Its reduction $D_e=M^e$ is nonempty and maximum. Give $D_e$ any
representation, extend it literally to $\pi_eM$, and use
\eqref{eq:maximum-cube-inward-lift} with sink bit $b$. The proof already
given covers both consecutive ranks and the full-cube identity case.
For every two-point fiber, the only outgoing $e$-label is on the
$(1-b)$-lift. Thus every $e$-edge points toward shore $b$.
For a rank-zero parent the unique map suffices, and any nonempty marked
section is that singleton.

Restrict this orientation to any marked section $A=M[\beta]$ in which
$e$ is not yet fixed. Coordinate-face restriction preserves ordinary
USOs, as proved above. Therefore it supplies a USO of $A$ with inward
boundary at $B=A[e=b]$, whenever $B$ is nonempty. This argument uses
only the actual parent $M$, even though $A$ need not itself be maximum.

We next verify the replacement statement needed for a noncube $B$.
Let an ample class $A$ have an ordinary USO with inward boundary at
$B=A[\gamma]\ne\varnothing$. Keep all exterior and boundary edges,
and replace the induced orientation on $B$ by any ordinary USO of $B$.
At a vertex of $B$, all outgoing directions remain internal and its
new outgoing cube is supplied by that USO. All outside outgoing cubes
are unchanged. A contained cube $L\subseteq A$ meeting $B$ intersects
it in a full face $L\cap B$ of $L$, even if $B$ is not a cube.
The old unique sink on that face was the unique sink of $L$, because
its boundary in $L$ is inward. Every outside word of $L$ still has
an outgoing edge within $L$ after replacement. The new face has one
sink, again a sink of $L$. Cubes disjoint from $B$ are unchanged.
This proves both ordinary-USO conditions with fixed exterior and boundary.

Refine the displayed chain by fixing one additional coordinate at a
time, retaining any constants and omitting duplicate fixed-coordinate
instructions. Every intermediate class is a nonempty marked section
of the same $M$. For each adjacent pair $A\supseteq B$, choose once
and for all the single-coordinate inward orientation of $A$ just
constructed. Start with the prescribed map on the final class $D$.
Working backwards, replace the map on $B$ in the chosen orientation
of $A$ by the map already obtained there. The replacement statement
proves validity at every step. At every word of an inner retained
class, its entire out-set stays the same when the next outer layer
is attached: the boundary contributes no outgoing direction there.
Hence all inward boundaries already obtained persist. All out-sets outside the final $D$ come from
the initially chosen layer exteriors, so do not depend on its map.
This proves the chain assertion and, by including $H$ in the chain,
the first assertion as well.
\end{proof}

\begin{corollary}[Ordinary gluing over a noncube coordinate face]
\label{cor:maxmin-noncube-face-gluing}
Let $S,F_1,\ldots,F_m$ be pairwise disjoint, and let nonempty ample
$A_i\subseteq\{0,1\}^{S\sqcup F_i}$ have the same nonempty section
$D=A_i\cap(\{0,1\}^S\times\{0\}^{F_i})$. Embed the pieces with
other exclusive blocks zero. Suppose $H=\bigcup_iA_i$ is ample.
If all but at most one piece are maximum-generated and the remaining
piece supplies any ordinary representation, then $H$ is represented,
preserving all out-sets of that piece. If all pieces are maximum-generated,
every ordinary USO on $D$ extends with inward boundary at $D$ and
with one fixed exterior. The intersection $D$ need not be a full cube.
\end{corollary}
\begin{proof}
The induced map of the exceptional piece on $D$ is an ordinary USO,
since $D$ is its coordinate-face intersection. Apply
Theorem~\ref{thm:maxmin-coordinate-face-inward} to choose each other
piece map inward at $D$ with this induced map. All orientations agree
on $D$, and at any word there outgoing exclusive directions belong
to at most the exceptional piece. Its outgoing cube supplies the
whole outgoing cube in the union. Outside $D$, all neighbors lie in
one piece, and every contained cube of $H$ belongs to a piece, by
the same exclusive-block argument as in
Theorem~\ref{thm:maxmin-common-cube-map-selection}. Thus the union
is an ordinary USO, with all exceptional out-sets fixed.
With no exception, apply the theorem to every piece and the same
arbitrary prescribed map on $D$. Fixed piece exteriors give a fixed
union exterior. Ampleness is a hypothesis here: for a noncube $D$
the common-shadow equality need not hold automatically.
\end{proof}

\begin{corollary}[Binary extentures and forbidden representations]
\label{cor:maxmin-binary-forbidden-representation}
Every hypothetical ample forbidden pc-minor of $\mathsf{MaxMin}$
that has a binary extenture is ordinarily representable. Consequently
any unrepresented ample forbidden MaxMin input shatters every pair of
its active coordinates.
An ample forbidden pc-minor of the dismantlable-ample class with a
binary extenture is likewise represented if at least one of the two
coordinate-face pieces of that cut is maximum-generated.
\end{corollary}
\begin{proof}
For a binary extenture $\{x_i=b,x_j=c\}$, put
$A=H[x_j=1-c]$ and $B=H[x_i=1-b]$. Then $H=A\cup B$ and
$D=A\cap B=H[x_i=1-b,x_j=1-c]$ is nonempty. Indeed the missing
pattern is minimal, so the projections realize the two patterns adjacent to the
missing pattern; a shortest path in the ample class between such
words must pass through the intersection, since the fourth pattern
is absent. Both pieces are proper coordinate-face minors.
They use the common block of all coordinates except $i,j$, with
one exclusive coordinate each, after translating the retained bits
to zero. The preceding corollary applies. For a forbidden MaxMin
input both pieces are MaxMin. For a forbidden dismantlable-ample
input both pieces are dismantlable and hence represented; the extra
hypothesis makes one piece an inward supplier.
\end{proof}

The replay extends prescribed acyclic and cyclic maps on the noncube
$129$-word section of the same $299$-word maximum parent. It also
retains an inward chain through a $66$-word noncube face, with identical
exteriors for the two target maps. A $195$-word amalgam over that face
preserves every out-set of its prescribed prism piece. Original parent
constants are retained. This does not give maximum ancestry of the
amalgam or maximum ancestry of arbitrary positive dismantlable pieces.
Simultaneous prescription on intersecting, nonnested faces is also
outside the theorem's chain assertion.
